\section{Supplied propensity scores and assignment information}

The calibrated procedure in \cref{obj:def:explicit-score-ledger} obtains assignment information from the original records. We now consider an experiment in which the entire true continuous propensity is supplied together with those records. This experiment supports two complementary conclusions: an attaining test over a broader class with continuous propensity and conditional mean functions, and an exact comparison of testing sensitivity on the original model.

A supplied-function test \leanref{sym:phior}{\ensuremath{\phi^{\mathrm e}}} is a jointly Borel map
\[
\phi^{\mathrm e}:
C([0,1])\times
\bigl([0,1]\times\{0,1\}\times\mathbb R\bigr)^n
\times[0,1]\longrightarrow[0,1].
\]
The function space carries the uniform topology. The remaining arguments are the independent original-record sample and an independent public uniform randomization variable. Under a law \(P\), the supplied function is \(e_P\); the formal risk definition below specifies the corresponding rejection expectation.

\subsection{The continuous-carrier experiment}

Let \(\mathcal V\) denote the admissible domain of public moment and smoothness tuples in \cref{obj:def:model}, and let \(\lambda\) be Lebesgue probability measure on \([0,1]\). The broader model retains uniform design, overlap, effect smoothness, the mean caps, and the conditional moment envelope. Its propensity and both conditional arm means have continuous representatives. The following definition specifies these restrictions and the maximum heterogeneity distance.

\begin{definitionv}{def:oracle-broad-model}[Continuous-carrier model \(\mathcal M_v^{\mathrm{e,br}}\)]
The class \(\mathcal M_v^{\mathrm{e,br}}\) consists of the original-record laws satisfying the six conditions below and admitting continuous representatives of their propensity and both conditional arm means. Here \(v=(p,\alpha,\beta,\gamma)\in\mathcal V\) is an admissible public tuple, \(\lambda\) is Lebesgue probability measure on \([0,1]\), and \(Q_{a,P}(dy\mid x)\) is the conditional outcome probability kernel in treatment arm \(a\).
\begin{itemize}
\item \textbf{(Uniform covariates.)} The covariate marginal satisfies
\[
 P_X=\lambda,
\]
as stipulated in \cref{obj:ass:uniform}.
\item \textbf{(Overlap.)} The continuous propensity representative satisfies
\[
 e_P([0,1])\subseteq[1/4,3/4],
\]
as stipulated in \cref{obj:ass:overlap}.
\item \textbf{(Effect smoothness.)} The continuous mean effect satisfies
\[
 \tau_P=m_{1,P}-m_{0,P}\in\mathcal H^\gamma(20),
\]
where \(\mathcal H^\gamma(20)\) consists of continuous functions on \([0,1]\) whose sup norm and \(\gamma\)-Hölder constant are at most \(20\), as stipulated in \cref{obj:ass:effect-smoothness}.
\item \textbf{(Baseline cap.)} The baseline mean satisfies
\[
 \|m_{0,P}\|_\infty\le1/2,
\]
as stipulated in \cref{obj:ass:baseline-cap}.
\item \textbf{(Effect cap.)} The mean effect satisfies
\[
 \|\tau_P\|_\infty\le1/2,
\]
as stipulated in \cref{obj:ass:effect-cap}.
\item \textbf{(Raw moment.)} The conditional outcome kernels satisfy
\[
 \max_{a\in\{0,1\}}\operatorname*{ess\,sup}_{x\sim\lambda}
 \int |y|^p Q_{a,P}(dy\mid x)\le10,
\]
as stipulated in \cref{obj:ass:raw-moment}.
\end{itemize}
For a function \(h\), write \([h]_\lambda\) for its equivalence class under Lebesgue almost-everywhere equality. Let \(\mathcal E_P\) denote the propensity equivalence class and let \(\mathcal G_{0,P},\mathcal G_{1,P}\) denote the conditional arm-mean equivalence classes. The continuous-representative requirement is
\[
 \exists f,g_0,g_1\in C([0,1]):\qquad
 \mathcal E_P=[f]_\lambda,\qquad
 \mathcal G_{0,P}=[g_0]_\lambda,\qquad
 \mathcal G_{1,P}=[g_1]_\lambda.
\]
Here \(C([0,1])\) is the space of continuous real functions on the unit interval with the uniform topology. These representatives are the existing functionals \(e_P,m_{0,P},m_{1,P}\). Membership depends on \(p\) and \(\gamma\), with the full tuple \(v\) retained as the public index.

Writing \(d(P)\) for the population Lebesgue \(L_2\) distance of \(\tau_P\) from its average, define
\[
 D_v^{\mathrm{e,br}}=\sup_{P\in\mathcal M_v^{\mathrm{e,br}}}d(P).
\]
\end{definitionv}

The continuous-carrier class \leanref{sym:MoracleBroad}{\ensuremath{\mathcal M_v^{\mathrm{e,br}}}} in \cref{obj:def:oracle-broad-model} permits continuous propensity and baseline mean functions with unrestricted moduli of continuity. Effect smoothness controls the approximation of heterogeneity, while overlap bounds the inverse-propensity weights. The conditional moment envelope controls outcome clipping uniformly across covariate values and treatment arms. The caps bound the scale of the conditional means and effect. Within this class, the null retains the unknown constant effect used in the original experiment.

\begin{definitionv}{def:oracle-broad-null}[Constant-effect null \(H_0^{\mathrm{e,br}}(v)\)]
The constant-effect null is
\[
 H_0^{\mathrm{e,br}}(v)=
 \left\{P\in\mathcal M_v^{\mathrm{e,br}}:
 \exists c\in[-1/2,1/2]\ \forall x\in[0,1],\
 \tau_P(x)=c\right\}.
\]
Membership in the continuous-carrier model imposes the restrictions in
\cref{obj:ass:uniform,obj:ass:overlap,obj:ass:effect-smoothness,obj:ass:baseline-cap,obj:ass:effect-cap,obj:ass:raw-moment}.
The unknown-constant effect condition is that of \cref{obj:ass:null-constancy}.
\end{definitionv}

Size is controlled uniformly over the constant-effect null \leanref{sym:HnulloracleBroad}{\ensuremath{H_0^{\mathrm{e,br}}(v)}} in \cref{obj:def:oracle-broad-null}. For a supplied-function test, write \leanref{sym:Riskor}{\ensuremath{\mathsf R_n^{\mathrm e}(P,\phi^{\mathrm e})}} for its rejection expectation when it receives \(e_P\). The following risk and radius measure the separation required to attain type-II error at most \(1/10\), subject to the same bound on size.

\begin{definitionv}{def:oracle-broad-risk}[Supplied-propensity risk \(B_n^{\mathrm{e,br}}(v,r)\)]
The supplied-propensity minimax type-II error is
\[
 B_n^{\mathrm{e,br}}(v,r)=
 \inf_{\phi^{\mathrm e}:\,
     \sup_{P\in H_0^{\mathrm{e,br}}(v)}
       \mathsf R_n^{\mathrm e}(P,\phi^{\mathrm e})\le1/10}
 \ \sup_{P\in\mathcal M_v^{\mathrm{e,br}}:\,d(P)\ge r}
       \left\{1-\mathsf R_n^{\mathrm e}(P,\phi^{\mathrm e})\right\},
\]
for parameters satisfying
\begin{itemize}
\item \textbf{(Sample size.)} \(n\ge2\).
\item \textbf{(Public tuple.)} \(v\in\mathcal V\), the admissible domain of moment and primitive smoothness tuples.
\item \textbf{(Separation.)} \(0<r<D_v^{\mathrm{e,br}}\), where \(D_v^{\mathrm{e,br}}\) is the supremum of the population effect distance over \(\mathcal M_v^{\mathrm{e,br}}\).
\end{itemize}
Here \(d(P)\) is the population Lebesgue \(L_2\) distance of the mean effect from its average. The admissible tests \(\phi^{\mathrm e}\) are precisely the jointly Borel maps taking a supplied continuous function, an original-record sample, and a public randomization variable to a rejection probability in \([0,1]\). Their rejection expectation is
\[
 \mathsf R_n^{\mathrm e}(P,\phi^{\mathrm e})
 =\int \phi^{\mathrm e}(e_P,\mathcal D_n,u)\,
       P^{\otimes n}(d\mathcal D_n)\,du.
\]
The sample \(\mathcal D_n\) consists of \(n\) independent original records with law \(P\); the integration in \(u\) is over an independent uniform variable on \([0,1]\). The supplied function \(e_P\) is the whole true continuous propensity.
\end{definitionv}

The broader-class risk \leanref{sym:BoracleBroad}{\ensuremath{B_n^{\mathrm{e,br}}(v,r)}} in \cref{obj:def:oracle-broad-risk} uses the original outcome records throughout. Its capped radius includes the model supremum as a saturation value, making the sensitivity measure well defined at every stipulated sample size.

\begin{definitionv}{def:oracle-broad-radius}[Capped radius \(r_n^{*,\mathrm{e,br}}(v)\)]
The capped supplied-propensity critical radius is
\[
 r_n^{*,\mathrm{e,br}}(v)=
 \inf\left(
 \left\{r\in(0,D_v^{\mathrm{e,br}}):
 B_n^{\mathrm{e,br}}(v,r)\le1/10\right\}
 \cup\left\{D_v^{\mathrm{e,br}}\right\}\right).
\]
Here \(D_v^{\mathrm{e,br}}\) is the supremum of the population effect distance over the continuous-carrier model and serves as the saturation value in the capped-radius convention.
\end{definitionv}

The capped radius \leanref{sym:rstaroracleBroad}{\ensuremath{r_n^{*,\mathrm{e,br}}(v)}} and its saturation value \leanref{sym:DoracleBroad}{\ensuremath{D_v^{\mathrm{e,br}}}} are defined in \cref{obj:def:oracle-broad-radius,obj:def:oracle-broad-model}. To attain the separation rate, the test uses clipped inverse-propensity outcomes and the histogram features of \cref{obj:def:projection}. Removing the constant histogram direction targets deviations from an unknown constant effect. The procedure then takes the inner product of two independent block averages.

\begin{algorithmv}{def:oracle-handle}[Supplied-propensity test \(\phi^{*,\mathrm e}_{n,v}\)]
Inputs are \(n\ge2\) original records, the public tuple \(v=(p,\alpha,\beta,\gamma)\), and a supplied continuous propensity \(e\).
\begin{enumerate}
\item Set the block size, moment exponent, separation exponent, and target scale to
\[
s=\lfloor n/2\rfloor,
\qquad
q=\frac{p-1}{p},
\qquad
E_0=\frac{2\gamma q}{2\gamma+q},
\qquad
a=s^{-E_0}.
\]
Use the evaluation blocks
\[
\mathcal I_1=\{1,\ldots,s\},
\qquad
\mathcal I_2=\{s+1,\ldots,2s\}.
\]
\item Select the clipping cutoff and histogram rank
\[
T=2^{\left\lceil\log_2(a^{-1/(p-1)})\right\rceil},
\qquad
J=2^{\left\lceil\log_2(a^{-1/\gamma})\right\rceil}.
\]
Define outcome clipping and the overlap-clipped propensity by
\[
\ell_T(y)=\max\{-T,\min\{y,T\}\},
\qquad
\widetilde e(x)=\max\{1/4,\min\{e(x),3/4\}\}.
\]
For every legal supplied true propensity, \(\widetilde e=e\).
\item Form the original-record inverse-propensity scores
\[
R_i=
\frac{A_i\ell_T(Y_i)}{\widetilde e(X_i)}
-
\frac{(1-A_i)\ell_T(Y_i)}{1-\widetilde e(X_i)}.
\]
Let \(F_J\) be the normalized histogram feature vector. Its constant direction and the centered block vectors are
\[
u=J^{-1/2}(1,\ldots,1)^T,
\qquad
V_b=\frac1s\sum_{i\in\mathcal I_b}
\bigl(F_J(X_i)-u\bigr)R_i,
\quad b\in\{1,2\}.
\]
These vectors represent the block scores obtained by applying the constant-centered histogram projection \(\Pi_J-1\).
\item Set the public clipping-bias and covariance bounds to
\[
b=20T^{1-p},
\qquad
\Lambda=\frac{80T^{2-p}}s.
\]
Output the rejection rule
\[
\phi^{*,\mathrm e}_{n,v}
=
\mathbf 1\!\left\{
\langle V_1,V_2\rangle
>
2b^2+1024\sqrt J\,\Lambda
\right\}.
\]
Its calibration and necessity properties are given by its singleton and canonical quadratic projections and the normalized common-propensity paired-tent converse.
\end{enumerate}
\end{algorithmv}

The supplied-propensity rule \leanref{sym:phistaror}{\ensuremath{\phi^{*,\mathrm e}_{n,v}}} in \cref{obj:def:oracle-handle} uses public clipping-bias and covariance bounds to set its rejection threshold. The tuning depends on the moment order and effect smoothness. The next result gives matching bounds for the broader-class critical radius, uniform calibration over its entire null, and power whenever the stated separation lies below the saturation value.



The conditional-mean identity in \cref{obj:thm:oracle-frontier-broad-class} explains the role of supplied assignment information: inverse weighting recovers the mean effect directly, and constant centering isolates its heterogeneity. Clipping controls the contribution of outcome tails, while effect smoothness controls histogram approximation. Their balance yields the supplied-propensity separation scale \leanref{sym:rhoor}{\ensuremath{\rho_n^{\mathrm e}(v)}} defined in that result. The matching lower bound uses null and alternative laws with the same constant propensity and close complete-record distributions. Thus the stated separation order remains necessary even when assignment information is identical across the comparison laws.

The radius bounds apply to every \(n\ge2\), with the cap accounting for finite-sample saturation. The procedure's power guarantee applies when its upper separation is strictly below the broader-model supremum. The public sample-size threshold in \cref{obj:thm:oracle-frontier-broad-class} ensures this condition by placing that separation at most halfway to the fixed witness distance. These statements distinguish an all-sample radius bound from a power guarantee over a nonempty range of alternatives.

\subsection{Comparison on the original model}

To measure the value of supplied assignment information on a fixed statistical class, we now use the original model of \cref{obj:def:model}. Its quantitative propensity and baseline smoothness restrictions are retained in both experiments being compared. The supplied-function test carrier remains the same, while the null and alternative suprema are taken over the original model. The following definitions make this shared-class comparison precise.

\begin{definitionv}{def:oracle-risk}[Supplied-propensity risk $B_n^{\mathrm e}(v,r)$]
\[
B_n^{\mathrm e}(v,r)=
\inf_{\phi^{\mathrm e}:\ \sup_{P\in H_0(v)}
\mathsf R_n^{\mathrm e}(P,\phi^{\mathrm e})\le 1/10}
\ \sup_{P\in\mathcal M_v:\ d(P)\ge r}
\bigl\{1-\mathsf R_n^{\mathrm e}(P,\phi^{\mathrm e})\bigr\}.
\]
The parameters satisfy:
\begin{itemize}
\item \textbf{(Sample size.)} \(n\ge 2\).
\item \textbf{(Public tuple.)} \(v\in\mathcal V\), the admissible domain of moment and primitive smoothness tuples.
\item \textbf{(Separation.)} \(0<r<D_v\), where \(D_v\) is the supremum of \(d(P)\) over \(\mathcal M_v\).
\end{itemize}
Here \(\mathcal M_v\) is the original-record model, \(H_0(v)\) consists of its laws with an unknown constant mean effect, and \(d(P)\) is the population Lebesgue \(L_2\) distance of the effect from its average. The infimum ranges over randomized rejection-probability maps \(\phi^{\mathrm e}\) receiving the whole true continuous propensity, and \(\mathsf R_n^{\mathrm e}(P,\phi^{\mathrm e})\) is their expected rejection probability with that propensity supplied.
\end{definitionv}

The shared-class risk \leanref{sym:Bor}{\ensuremath{B_n^{\mathrm e}(v,r)}} in \cref{obj:def:oracle-risk} uses the original-model null for its size constraint. Its radius therefore uses the original-model supremum as the saturation value.

\begin{definitionv}{def:oracle-radius}[Supplied-propensity radius $r_n^{*,\mathrm e}(v)$]
\[
r_n^{*,\mathrm e}(v)=
\inf\Bigl(
\bigl\{r\in(0,D_v):B_n^{\mathrm e}(v,r)\le 1/10\bigr\}
\cup\{D_v\}
\Bigr).
\]
Here \(B_n^{\mathrm e}(v,r)\) is the supplied-propensity minimax type-II error on the original model under the size constraint \(1/10\), and \(D_v\) is the supremum of the population effect distance over that model.
\end{definitionv}

The shared-class capped radius \leanref{sym:rstaror}{\ensuremath{r_n^{*,\mathrm e}(v)}} of \cref{obj:def:oracle-radius} is compatible with the broader-class attainment result. The original model is contained in the continuous-carrier class, and the finite lower-bound witnesses in \cref{obj:thm:oracle-frontier-broad-class} belong to both classes. The next result records the matching separation bounds and attainment on the original model, providing the supplied-propensity benchmark for the radius comparison.



The compatibility in \cref{obj:thm:oracle-frontier-resolved} rests on both attainment and necessity. The supplied-propensity procedure controls size and power on the original class, and common-propensity witnesses establish the corresponding lower separation bound within that class. Consequently, the original-model comparison uses matched statistical classes in its numerator and denominator. Its finite-sample power condition is expressed relative to the original-model supremum; the broader-class result uses its own supremum.

We measure preservation of separation order by a uniformly bounded ratio of the unknown- and supplied-propensity capped radii. The following region includes every admissible tuple for which that ratio remains bounded across sample sizes.

\begin{definitionv}{def:preservation-region}[Propensity preservation region $\mathcal R$]
\[
\mathcal R=
\left\{
v\in\mathcal V:
\sup_{n\ge 2}
\frac{r_n^*(v)}{r_n^{*,\mathrm e}(v)}
<\infty
\right\}.
\]
Here \(\mathcal V\) is the admissible domain of public moment and primitive smoothness tuples. The radii \(r_n^*(v)\) and \(r_n^{*,\mathrm e}(v)\) are the capped critical separation radii on the same original-record model with unknown and supplied true continuous propensity, respectively.
\end{definitionv}

The preservation region \leanref{sym:Region}{\ensuremath{\mathcal R}} in \cref{obj:def:preservation-region} follows from comparing the matched rates in \cref{obj:thm:heavy-tail-comparison,obj:thm:oracle-frontier-resolved}. The scalar comparison quantity \(F(p)\) defined in \cref{obj:def:sharp-frontier-scales} determines whether the separation exponents coincide. The next result gives the exact boundary, including its equality case and bounds for every capped finite-sample radius ratio.



The boundary in \cref{obj:thm:exact-propensity-preservation-boundary} expresses the value of assignment information through the difference between the two separation exponents. When \(F(p)\ge1\), the unknown-propensity experiment attains the supplied-propensity separation order on the same original class. At \(F(p)=1\), the ratio remains between positive, sample-size-independent constants. When \(F(p)<1\), supplying the true propensity improves sensitivity by a polynomial factor whose exponent is stated explicitly in the result.

These comparisons concern orders of critical separation, with finite-sample constants and saturation incorporated into the capped radii. They hold for every \(n\ge2\), and their constants admit uniform bounds on compact subsets of the admissible parameter domain. Together with the broader-class attainment result, they distinguish testing with supplied assignment information from the exact region where the original unknown-propensity experiment attains the same separation order. Requirements for a particular estimated-propensity procedure are outside this experiment comparison.
